% Provenance (never printed): sources/book2/BRIEFS.md, g12-06, with the sync ruling: no WGI (licence); the opening chart is a
% Maddison pair different from Book 1 g9-06 (Haiti and the Dominican Republic). Ideas are nonrival; incentives to invent;
% creative destruction; institutions recalled; diffusion; growth and the environment. Real numbers:
% figdata/grade-12/06-ideas-institutions-growth/ (Maddison HTI, DOM 1950-2022; WDI energy intensity and world GDP 2000-2021).
% Ledger: Bloom et al. 2020 abstract, Longitude rewards, BLS Business Employment Dynamics 2026 Q1. Patent is Book 1's
% (recalled). The Korean photograph is Book 1's and is not repeated; no governance index is printed. The medicine of the
% weekend problem is invented.
\chapter{Ideas, Institutions and Growth}\label{ch:g12:ideas-institutions-growth}

Haiti and the Dominican Republic share one island, its climate and its storms. In 1950 output per person was about the same in both: 1790 dollars in Haiti and 1637 in the
Dominican Republic, in the Maddison Project's estimates. By 2022 the Dominican figure was 17\,714 dollars. Haiti's was 1546, lower than in 1950.
% figdata: island (1950 HTI 1790, DOM 1637; 2022 HTI 1546, DOM 17714)

The Solow model says that lasting growth needs better technology, and that capital follows technology. It does not say where better technology comes from, nor why some
countries use the best methods and others do not. This chapter takes up both questions: \omterm{def:g12:ideas-institutions-growth:idea}{ideas}, and the rules under which people produce and use them.

\begin{openingchart}{GDP per person in Haiti and in the Dominican Republic, 1950--2022, in 2011 dollars, logarithmic scale.}{Maddison Project Database}
\begin{tikzpicture}
\begin{axis}[omtime, width=0.84\linewidth, height=5.4cm, xmin=1950, xmax=2022, ymode=log, ymin=1000, ymax=25000, log ticks with fixed point,
  ytick={1000,2000,5000,10000,20000}, xtick={1950,1960,...,2020}, ylabel={dollars of 2011}]
\addplot[omDef, thick] table[x=year, y=DOM] {figdata/grade-12/06-ideas-institutions-growth/out/island.csv};
\addplot[omThm, thick] table[x=year, y=HTI] {figdata/grade-12/06-ideas-institutions-growth/out/island.csv};
\node[omDef, font=\scriptsize, anchor=south] at (axis cs:2000,8500) {Dominican Republic};
\node[omThm, font=\scriptsize, anchor=north] at (axis cs:2000,1400) {Haiti};
\end{axis}
\end{tikzpicture}
\end{openingchart}

\begin{recall}
From the grades 6--9 volume: patents; institutions, natural experiments and the comparison of planned and market economies. From \cref{ch:g11:externalities}: external
benefits. From \cref{ch:g11:public-goods-and-commons}: rival and \omterm{def:g11:public-goods-and-commons:rival}{excludable goods}. From \cref{ch:g12:growth-and-production-function} and \cref{ch:g12:solow-first-look}: \omterm{def:g12:growth-and-production-function:tfp}{total
factor productivity}, and why growth in the Solow model needs it to rise.
\end{recall}

\section{Ideas are nonrival}\label{sec:g12:ideas-institutions-growth:ideas}

\begin{definition}[Idea]\label{def:g12:ideas-institutions-growth:idea}
An \emph{idea}\index{idea}, in economics, is a set of instructions for making or doing something: a design, a formula, a method, a program. It is a nonrival good: one
person's use of it does not leave less for anyone else.
\end{definition}

A machine can be used by one worker at a time; the design of the machine can be used by every factory at once. That is why \omterm{def:g12:ideas-institutions-growth:idea}{ideas} can sustain growth where capital cannot:
an \omterm{def:g12:ideas-institutions-growth:idea}{idea}'s returns do not diminish as more people use it. But a nonrival good that others can copy brings its inventor little, and so may not be produced at all: an \omterm{def:g12:ideas-institutions-growth:idea}{idea}
has large external benefits.

\section{The incentives to invent}\label{sec:g12:ideas-institutions-growth:incentives}

\begin{definition}[Research and development]\label{def:g12:ideas-institutions-growth:rnd}
\emph{Research and development}\index{research and development} is work done to produce new \omterm{def:g12:ideas-institutions-growth:idea}{ideas}: new products, new processes, or new knowledge that may lead to them.
\end{definition}

Three ways to pay for it are in use. A patent gives the inventor the sole right to use an \omterm{def:g12:ideas-institutions-growth:idea}{idea} for a period, often twenty years: the inventor charges a monopoly price, and
the \omterm{def:g10:what-prices-do:dwl}{deadweight loss} of monopoly is the cost of the incentive. Public funding of research, in universities and laboratories, pays for \omterm{def:g12:ideas-institutions-growth:idea}{ideas} whose results are given away.
And a prize pays for a result defined in advance.

\begin{definition}[Innovation prize]\label{def:g12:ideas-institutions-growth:prize}
An \emph{innovation prize}\index{innovation prize} is a sum promised in advance to whoever first achieves a specified result, after which the \omterm{def:g12:ideas-institutions-growth:idea}{idea} can be used freely.
\end{definition}

\begin{omfigure}
\begin{tikzpicture}
\begin{axis}[width=0.62\linewidth, height=5.4cm, xmin=0, xmax=105, ymin=0, ymax=105, xlabel={quantity}, ylabel={price},
  xtick={40,80}, ytick={20,60,100}]
\fill[omThm!25] (axis cs:0,60) -- (axis cs:40,60) -- (axis cs:40,20) -- (axis cs:0,20) -- cycle;
\fill[omProp!35] (axis cs:40,60) -- (axis cs:80,20) -- (axis cs:40,20) -- cycle;
\fill[omDef!15] (axis cs:0,100) -- (axis cs:40,60) -- (axis cs:0,60) -- cycle;
\addplot[omDef, thick, domain=0:100, samples=2] {100 - x} node[pos=0.12, above right, font=\scriptsize] {demand};
\addplot[black, thick, domain=0:105, samples=2] {20} node[pos=0.95, above, font=\scriptsize] {cost};
\draw[dashed, gray] (axis cs:40,0) -- (axis cs:40,60);
\node[font=\scriptsize] at (axis cs:20,40) {inventor's profit};
\node[font=\scriptsize, anchor=west] at (axis cs:47,32) {lost};
\node[font=\scriptsize] at (axis cs:12,70) {buyers};
\end{axis}
\end{tikzpicture}

{\small The patent's trade-off, with the numbers of the weekend problem. While the patent lasts, the inventor charges 60 instead of the cost of 20: the profit (rectangle) is
the reward for inventing, and the triangle marked ``lost'' is the \omterm{def:g10:what-prices-do:dwl}{deadweight loss}, the treatments not sold that buyers valued above their cost. A prize paid instead lets
the medicine sell at cost from the start.}
\end{omfigure}

\begin{record}[The longitude prizes]
% ledger: g12:longitude
A ship at sea could not tell how far east or west it was. In 1714 the British Parliament offered 10\,000 pounds for a method accurate to within one degree of longitude,
15\,000 for 40 minutes and 20\,000 for 30 minutes. A clockmaker received a first 250 pounds in 1737 to improve his sea clock; after decades of trials and
disputes over whether his result met the terms, Parliament granted him 8750 pounds in 1773, bringing his total to 23\,065 pounds over thirty-six years. The record shows
both the strength of a prize, a result anyone may then use, and its weakness: someone must judge whether the result has been reached.
\end{record}

\begin{definition}[Creative destruction]\label{def:g12:ideas-institutions-growth:creative-destruction}
\emph{Creative destruction}\index{creative destruction} is the process by which new products, methods and firms replace old ones, so that growth comes with the closing of
firms and the loss of jobs as well as with new ones.
\end{definition}

% ledger: g12:bed-2026q1
In the United States in the first quarter of 2026, expanding and opening private businesses added 7.4 million jobs, and shrinking and closing ones lost 7.3 million: each
about 5.5\,\% of private employment, for a net change of 88\,000. In that quarter 323\,000 new establishments opened. Behind a small net change in employment lies a large
reallocation of workers from some firms to others.

\begin{model}{Ideas-driven growth (toy)}
\assumes \omterm{def:g12:growth-and-production-function:tfp}{Total factor productivity} grows in proportion to the number of researchers: $g_A = \theta R$, with $\theta$ constant.
\predicts More researchers bring faster growth, and a growing number of researchers brings ever faster growth.
\tested Research effort has risen substantially for decades while productivity growth has not: the number of researchers needed to double the density of computer chips is
  more than 18 times larger than in the early 1970s (Bloom, Jones, Van Reenen and Webb 2020).
\breaks \omterm{def:g12:ideas-institutions-growth:idea}{Ideas} are getting harder to find: $\theta$ falls as the easy \omterm{def:g12:ideas-institutions-growth:idea}{ideas} are used up, and the same growth needs ever more research.
\end{model}
% ledger: g12:bloom2020

\section{Institutions}\label{sec:g12:ideas-institutions-growth:institutions}

The grades 6--9 volume compared countries split by history, the two Koreas, the two Germanys, and China before and after 1978, and found that the rules under which people
own, trade and decide made great differences to income. The same goes for \omterm{def:g12:ideas-institutions-growth:idea}{ideas}: a firm invests in research, and a worker in skills, only if they expect to keep a good share
of the return; a patent is worth little where courts do not enforce it. The opening chart is consistent with that, but a pair of countries is not a natural experiment:
Haiti and the Dominican Republic differ in their history, their politics and the disasters that struck them, and the chart alone cannot say which difference mattered most.

\section{Why the best methods do not spread everywhere}\label{sec:g12:ideas-institutions-growth:diffusion}

\begin{definition}[Technology diffusion]\label{def:g12:ideas-institutions-growth:diffusion}
\emph{Technology diffusion}\index{technology diffusion} is the spread of a new method or product from those who first use it to other firms and countries.
\end{definition}

\begin{omfigure}
\begin{tikzpicture}
\begin{axis}[width=0.66\linewidth, height=5.0cm, xmin=0, xmax=70, ymin=0, ymax=1.05, xlabel={years after invention}, ylabel={share of users},
  ytick={0,0.5,1}, grid=major, grid style={gray!20}]
\addplot[omDef, thick, domain=0:70, samples=100] {1/(1+exp(-0.3*(x-20)))};
\node[omDef, font=\scriptsize, anchor=east] at (axis cs:23,0.82) {first adopters};
\addplot[omThm, thick, domain=0:70, samples=100] {1/(1+exp(-0.3*(x-45)))};
\node[omThm, font=\scriptsize, anchor=west] at (axis cs:50,0.3) {late adopters};
\draw[<->, gray] (axis cs:20,0.5) -- node[below, font=\scriptsize, black] {lag} (axis cs:45,0.5);
\end{axis}
\end{tikzpicture}

{\small Diffusion as an S-curve (logistic, admitted): use spreads slowly at first, fast once many have adopted, and slowly again near saturation. A country that starts
later follows the same curve with a lag; a lag, not the shape, can make much of the difference between rich and poor countries.}
\end{omfigure}

New methods spread slowly: firms must learn them, buy the machines that carry them, train workers and change how they are organised, and a firm sheltered from competition
has less reason to bother. Where firms can enter freely and must compete, the best methods spread faster.

\section{Growth and the environment}\label{sec:g12:ideas-institutions-growth:environment}

\begin{definition}[Energy intensity]\label{def:g12:ideas-institutions-growth:energy-intensity}
\emph{Energy intensity}\index{energy intensity} is the primary energy used per unit of GDP.
\end{definition}

\begin{omchart}{Energy intensity of GDP, megajoules of primary energy per dollar of GDP at constant purchasing power parity, five countries and the world, 2000--2021.}{World
Bank, World Development Indicators (EG.EGY.PRIM.PP.KD)}
\begin{tikzpicture}
\begin{axis}[omtime, width=0.84\linewidth, height=5.4cm, xmin=2000, xmax=2021, ymin=0, ymax=12, xtick={2000,2005,...,2020}, ylabel={megajoules per dollar}]
\addplot[omDef, thick] table[x=year, y=USA] {figdata/grade-12/06-ideas-institutions-growth/out/intensity.csv};
\addplot[omThm, thick] table[x=year, y=CHN] {figdata/grade-12/06-ideas-institutions-growth/out/intensity.csv};
\addplot[omProp, thick] table[x=year, y=IND] {figdata/grade-12/06-ideas-institutions-growth/out/intensity.csv};
\addplot[omMeth, thick] table[x=year, y=DEU] {figdata/grade-12/06-ideas-institutions-growth/out/intensity.csv};
\addplot[gray, thick] table[x=year, y=GBR] {figdata/grade-12/06-ideas-institutions-growth/out/intensity.csv};
\addplot[black, thick, dashed] table[x=year, y=WLD] {figdata/grade-12/06-ideas-institutions-growth/out/intensity.csv};
\node[omThm, font=\scriptsize, anchor=south west] at (axis cs:2001,10.9) {China};
\node[omProp, font=\scriptsize, anchor=south west] at (axis cs:2001,7.2) {India};
\draw[omDef] (axis cs:2009,5.5) -- (axis cs:2009,4.45);
\node[omDef, font=\scriptsize, anchor=north] at (axis cs:2009,4.45) {United States};
\node[black, font=\scriptsize, anchor=south] at (axis cs:2016,5.0) {world};
\node[omMeth, font=\scriptsize, anchor=north] at (axis cs:2008,3.3) {Germany};
\node[gray, font=\scriptsize, anchor=north] at (axis cs:2018,2.3) {United Kingdom};
\end{axis}
\end{tikzpicture}
\end{omchart}
% figdata: intensity (USA 6.72 -> 4.24; CHN 10.91 -> 6.30; WLD 6.05 -> 4.54); intensity_key (gdp_ratio 2.0, energy_ratio 1.5)

Growth uses resources and produces pollution, and innovation has changed how much. Each dollar of world output used a quarter less primary energy in 2021 than in 2000,
from 6.05 to 4.54 megajoules; in the United States \omterm{def:g12:ideas-institutions-growth:energy-intensity}{energy intensity} fell by about 2.2\,\% a year. Innovation, and a shift of output towards services, have done that. What they have not done is cut the total:
world output doubled over those years, so world energy use still rose by half. Whether output can keep growing while the use of energy from fossil fuels, and the emissions
that come with it, fall is a question about the speed of innovation and about the prices that steer it: the \omterm{def:g11:externalities:private-cost}{external costs} of \cref{ch:g11:externalities}.

\section{Exercises}

\begin{exercise}[$\star$]\label{exo:g12:ideas-institutions-growth:1}
Rival or nonrival: a recipe; a cake; a computer program; a delivery truck?
\end{exercise}

\begin{exercise}[$\star$]\label{exo:g12:ideas-institutions-growth:2}
Rival or nonrival: a mathematical theorem; a seat on a train; a song; a dose of vaccine; the formula of the vaccine?
\end{exercise}

\begin{exercise}[$\star$]\label{exo:g12:ideas-institutions-growth:3}
From the opening chart, give both countries' GDP per person in 1950 and in 2022, and the ratio in 2022.
\end{exercise}

\begin{exercise}[$\star$]\label{exo:g12:ideas-institutions-growth:4}
From the opening chart, in which decade did the two countries begin to diverge? What was the Dominican figure in 1980 and the Haitian one?
\end{exercise}

\begin{exercise}[$\star$]\label{exo:g12:ideas-institutions-growth:5}
Demand is $P = 50 - Q$ and the \omterm{def:g11:marginal-reasoning:marginal-benefit}{marginal cost} is 10. A patent holder sells at the monopoly price. Find price, quantity, profit and the \omterm{def:g10:what-prices-do:dwl}{deadweight loss}.
\end{exercise}

\begin{exercise}[$\star\star$]\label{exo:g12:ideas-institutions-growth:6}
A patent brings a profit of 100 a year for 20 years. Find its \omterm{def:g12:present-value:present-value}{present value} at 4\,\%.
\end{exercise}

\begin{exercise}[$\star\star$]\label{exo:g12:ideas-institutions-growth:7}
An inventor can choose a prize of 1000 now or a patent bringing 80 a year for 20 years. Which is worth more at 5\,\%?
\end{exercise}

\begin{exercise}[$\star\star$]\label{exo:g12:ideas-institutions-growth:8}
In the weekend problem's market, the patent is lengthened from 20 to 25 years. By how much does its \omterm{def:g12:present-value:present-value}{present value} at 5\,\% rise? What does society give up in exchange?
\end{exercise}

\begin{exercise}[$\star\star$]\label{exo:g12:ideas-institutions-growth:9}
American \omterm{def:g12:ideas-institutions-growth:energy-intensity}{energy intensity} fell from 6.72 megajoules per dollar in 2000 to 4.24 in 2021. Find the average annual rate of change.
\end{exercise}

\begin{exercise}[$\star\star$]\label{exo:g12:ideas-institutions-growth:10}
From the energy-intensity chart, give China's value in 2000 and 2021. By what fraction did it fall? Compare with the world.
\end{exercise}

\begin{exercise}[$\star\star$]\label{exo:g12:ideas-institutions-growth:11}
The share of firms using a method is $1/(1 + e^{-0.3(t - 20)})$, $t$ in years. When does it reach 10\,\%, 50\,\% and 90\,\%? A second country follows the same curve 25 years
later: when does it reach 50\,\%?
\end{exercise}

\begin{exercise}[$\star\star\star$]\label{exo:g12:ideas-institutions-growth:12}
What evidence contradicts the toy model of ideas-driven growth? Use the model box, and say which assumption fails.
\end{exercise}

\begin{exercise}[$\star\star\star$]\label{exo:g12:ideas-institutions-growth:13}
Does the opening chart show that institutions cause income differences? Name two other explanations, and the kind of evidence that would separate them.
\end{exercise}

\begin{exercise}[$\star\star\star$]\label{exo:g12:ideas-institutions-growth:14}
In the toy model, let $g_A = 0.1R$, with $g_A$ in per cent a year and $R$ the number of researchers per 1000 workers. Find $g_A$ for $R = 10$ and $R = 20$. If $\theta$
halves while $R$ doubles, what happens to growth, and which record of the chapter does that resemble?
\end{exercise}

\begin{exercise}[$\star\star\star$]\label{exo:g12:ideas-institutions-growth:15}
Who loses from \omterm{def:g12:ideas-institutions-growth:creative-destruction}{creative destruction}, and who gains? Use the job-flow figures of the chapter, and link the answer to the measures of inequality of \cref{ch:g11:inequality}.
\end{exercise}

\section{Problem: Two Inventors in Brenland}

\begin{problem}[{Weekend problem --- a patent or a prize}]
\label{pb:g12:ideas-institutions-growth:1}
In a town in Brenland, two inventors race to develop a medicine. Yearly demand is $P = 100 - Q$ ($Q$ in thousands of treatments, $P$ in dalens), and each treatment costs 20
to make. The winner can take a 20-year patent, or the government can offer a prize instead and let anyone make the medicine. Use a \omterm{def:g12:present-value:present-value}{discount rate} of 5\,\%.

\medskip
\noindent\textbf{Part I --- The patent.}
\begin{enumerate}
\item Find the monopoly quantity and price.
\item Find the yearly profit.
\item Find the yearly \omterm{def:g10:surplus-and-efficiency:consumer-surplus}{consumer surplus} and \omterm{def:g10:what-prices-do:dwl}{deadweight loss} under the patent.
\item Find the \omterm{def:g12:present-value:present-value}{present value} of 20 years of profit.
\item Once the patent ends, what are the price, quantity and \omterm{def:g10:surplus-and-efficiency:consumer-surplus}{consumer surplus}?
\end{enumerate}

\medskip
\noindent\textbf{Part II --- The prize.}
\begin{enumerate}[resume]
\item What prize, paid now, gives the inventor the same \omterm{def:g12:present-value:present-value}{present value} as the patent?
\item With the prize, at what price is the medicine sold from the start?
\item Find the \omterm{def:g12:present-value:present-value}{present value} of the \omterm{def:g10:what-prices-do:dwl}{deadweight loss} avoided over 20 years.
\item Find the \omterm{def:g12:present-value:present-value}{present value} of the extra \omterm{def:g10:surplus-and-efficiency:consumer-surplus}{consumer surplus} over 20 years, and compare it with the prize.
\item Development costs 15\,000 thousand dalens. Is the medicine developed under the patent? Under the prize?
\end{enumerate}

\medskip
\noindent\textbf{Part III --- What can go wrong.}
\begin{enumerate}[resume]
\item The government underestimates demand and offers 12\,000. What happens?
\item It overestimates and offers 30\,000. Who gains, and at whose expense?
\item Why does a prize require the government to know more than a patent does?
\item Both inventors spend 15\,000 racing, and only one wins. What is the \omterm{def:g11:externalities:private-cost}{social cost} of the race? Is it wasted?
\item The longitude prizes took decades to pay out. What does that record suggest about defining the result in advance?
\end{enumerate}

\medskip
\noindent\textbf{Part IV --- Longer protection.}
\begin{enumerate}[resume]
\item The patent is lengthened to 25 years. Find its new \omterm{def:g12:present-value:present-value}{present value}.
\item With development costing 21\,000, which length of patent leads to the medicine being developed?
\item What is the cost to buyers of the five extra years, in \omterm{def:g12:present-value:present-value}{present value} of \omterm{def:g10:what-prices-do:dwl}{deadweight loss}?
\item Why might a longer patent bring more medicines but each at a higher cost to buyers?
\item State the result of the problem: the prize that gives the inventor the same \omterm{def:g12:present-value:present-value}{present value} as the patent, and the \omterm{def:g10:what-prices-do:dwl}{deadweight loss} it avoids.
\end{enumerate}
\end{problem}
